\documentclass[11pt,a4paper]{article}
\usepackage[utf8]{inputenc}
\usepackage[T1]{fontenc}
\usepackage[spanish]{babel}
\usepackage{graphicx}
\usepackage{float}
\usepackage[margin=2.2cm]{geometry}
\usepackage{enumitem}
\usepackage[hidelinks]{hyperref}

\title{Guía de observación\\EDA 02: revisión dirigida de hallazgos}
\author{Proyecto AgroCebada}
\date{}

\newcommand{\edafigure}[3]{
  \subsection{#1}
  \begin{figure}[H]
    \centering
    \includegraphics[width=0.96\textwidth,height=0.68\textheight,keepaspectratio]{\detokenize{#2}}
    \caption{#1}
  \end{figure}
  \textbf{Preguntas guía}
  #3
}

\begin{document}
\maketitle
\tableofcontents
\newpage

\section{Propósito}
Este segundo EDA profundiza en los principales hallazgos del análisis inicial:
calidad de la variable de nubosidad, heterogeneidad de cobertura Landsat,
comparación exacta Planet--Sentinel-2 y saltos temporales de NDVI. Para hacer
comparable la estacionalidad de nubosidad, la revisión mensual se limita a 2025.

\section{Nubosidad en 2025}

\edafigure{Nubosidad mensual por sensor en 2025}
{../../../outputs/eda/01_tablas_maestras/02_revision_dirigida/01_nubosidad_por_mes_2025.png}{
\begin{itemize}[nosep]
  \item ¿Qué meses presentan mayor nubosidad mediana y dispersión?
  \item ¿Los tres sensores muestran el mismo patrón estacional?
  \item ¿Las diferencias pueden reflejar reglas de enmascaramiento distintas?
\end{itemize}}

\edafigure{Categorías mensuales de nubosidad: Sentinel-2}
{../../../outputs/eda/01_tablas_maestras/02_revision_dirigida/02_nubosidad_categorias_sentinel_2_2025.png}{
\begin{itemize}[nosep]
  \item ¿Qué proporción mensual se concentra exactamente en 0 o 100?
  \item ¿Los valores intermedios varían con la temporada?
  \item ¿La concentración en extremos sugiere revisar cómo se calculó la variable?
\end{itemize}}

\edafigure{Categorías mensuales de nubosidad: Landsat}
{../../../outputs/eda/01_tablas_maestras/02_revision_dirigida/02_nubosidad_categorias_landsat_2025.png}{
\begin{itemize}[nosep]
  \item ¿En qué meses aparecen datos faltantes de nubosidad?
  \item ¿Los valores 0 y 100 dominan de manera similar a Sentinel-2?
  \item ¿La completitud cambia junto con la temporada de lluvias?
\end{itemize}}

\edafigure{Categorías mensuales de nubosidad: Planet}
{../../../outputs/eda/01_tablas_maestras/02_revision_dirigida/02_nubosidad_categorias_planet_2025.png}{
\begin{itemize}[nosep]
  \item ¿La composición mensual es estable durante 2025?
  \item ¿Qué meses tienen más observaciones clasificadas con 100 por ciento?
  \item ¿Las categorías son directamente comparables con los otros sensores?
\end{itemize}}

\edafigure{Completitud frente a nubosidad mediana}
{../../../outputs/eda/01_tablas_maestras/02_revision_dirigida/03_completitud_vs_nubosidad.png}{
\begin{itemize}[nosep]
  \item ¿Los meses más nublados tienen menor porcentaje de datos válidos?
  \item ¿Landsat se separa de Sentinel-2 y Planet en completitud?
  \item ¿Hay suficientes meses para sostener una relación consistente?
\end{itemize}}

\section{Heterogeneidad de cobertura}

\edafigure{Cobertura Landsat ordenada por parcela}
{../../../outputs/eda/01_tablas_maestras/02_revision_dirigida/04_cobertura_landsat_ordenada.png}{
\begin{itemize}[nosep]
  \item ¿La cobertura forma grupos claramente diferenciados?
  \item ¿Cuántas parcelas quedan muy por debajo o por encima de la mediana?
  \item ¿Debe estratificarse la interpretación posterior según nivel de cobertura?
\end{itemize}}

\section{Planet frente a Sentinel-2 en misma parcela y fecha}

\edafigure{Comparación de NDVI}
{../../../outputs/eda/01_tablas_maestras/02_revision_dirigida/05_comparacion_planet_sentinel_ndvi_promedio.png}{
\begin{itemize}[nosep]
  \item ¿Qué tan cerca están los puntos de la línea de igualdad?
  \item ¿Se observa un sesgo constante o dependiente del nivel de NDVI?
  \item ¿Hay valores extremos que deban rastrearse a parcelas y fechas concretas?
\end{itemize}}

\edafigure{Comparación de EVI}
{../../../outputs/eda/01_tablas_maestras/02_revision_dirigida/05_comparacion_planet_sentinel_evi_promedio.png}{
\begin{itemize}[nosep]
  \item ¿La concordancia visual es mejor o peor que para NDVI?
  \item ¿La nube de puntos es simétrica alrededor de la línea de igualdad?
  \item ¿Existen rangos donde un sensor responde con mayor sensibilidad?
\end{itemize}}

\edafigure{Comparación de LAI}
{../../../outputs/eda/01_tablas_maestras/02_revision_dirigida/05_comparacion_planet_sentinel_lai_promedio.png}{
\begin{itemize}[nosep]
  \item ¿Qué magnitud tiene el desplazamiento respecto a la igualdad?
  \item ¿El sesgo parece aditivo, proporcional o no lineal?
  \item ¿Por qué una correlación alta no basta para considerar equivalentes ambos LAI?
\end{itemize}}

\section{Saltos temporales de NDVI}

\edafigure{Cambios abruptos de NDVI frente a nubosidad en Sentinel-2}
{../../../outputs/eda/01_tablas_maestras/02_revision_dirigida/06_saltos_ndvi_vs_nubosidad.png}{
\begin{itemize}[nosep]
  \item ¿Los cambios absolutos grandes se concentran con nubosidad elevada?
  \item ¿También existen saltos fuertes con nubosidad registrada como cero?
  \item ¿Los saltos pueden deberse al intervalo temporal, al ciclo agrícola o a problemas de calidad?
  \item ¿Qué casos del CSV de extremos deben inspeccionarse individualmente?
\end{itemize}}

\section{Respuestas y hallazgos}

\subsection{Nubosidad en 2025}
\begin{enumerate}
  \item \textbf{Patrón mensual.} Junio es el mes con mayor nubosidad media en
  Sentinel-2 (66.0\%) y Landsat (74.4\%); Planet alcanza su máximo en septiembre
  (56.8\%). Los mínimos medios son noviembre para Sentinel-2 (1.8\%), abril para
  Landsat (0.3\%) y marzo para Planet (6.5\%). Hay estacionalidad, pero no es
  idéntica entre sensores y puede mezclar clima con diferencias de adquisición
  y enmascaramiento.

  \item \textbf{Sentinel-2.} En los meses de 2025, la categoría 0 representa
  entre 23.7\% y 95.4\%, la intermedia entre 3.4\% y 27.3\%, y 100 entre 1.1\%
  y 60.0\%. Junio tiene mediana de 100. La fuerte variación estacional y la
  concentración en extremos justifican revisar la definición de nubosidad.

  \item \textbf{Landsat.} El valor 0 fluctúa entre 21.0\% y 93.7\%, los valores
  intermedios entre 0.6\% y 22.5\%, y 100 entre 0.1\% y 68.1\%. Además, en todos
  los meses hay faltantes, entre 4.9\% y 10.0\% en 2025. Junio y septiembre
  tienen mediana de 100. Landsat es el sensor con mayor proporción global de
  observaciones en 100 y el único con datos faltantes de nubosidad.

  \item \textbf{Planet.} El valor 0 representa entre 31.1\% y 90.5\% mensual,
  la categoría intermedia entre 3.9\% y 23.0\%, y 100 entre 4.4\% y 46.0\%.
  Ningún mes de 2025 tiene mediana de 100 ni faltantes en esta variable. Su
  máximo mensual ocurre más tarde que en Sentinel-2 y Landsat.

  \item \textbf{Completitud.} Sentinel-2 y Planet tienen 100\% de completitud
  mensual. Landsat varía entre 87.47\% y 97.46\%, con media de 93.53\%. La menor
  completitud Landsat ocurre en enero de 2023, cuya mediana de nubosidad es 0;
  por tanto, no hay evidencia de que la falta de nubosidad aumente simplemente
  cuando la nubosidad mediana es alta.
\end{enumerate}

\subsection{Cobertura Landsat}
\begin{enumerate}[resume]
  \item Las 197 parcelas tienen el mismo rango general, del 5 de enero de 2022
  al 31 de diciembre de 2025, pero no el mismo número de observaciones. La
  cobertura va de 195 a 391, con mediana de 211, media de 246.8 y desviación
  estándar de 64.8. El cuartil superior comienza en 272, mostrando un grupo con
  cobertura claramente mayor. Conviene reportar resultados por nivel de
  cobertura y evitar que las parcelas con más registros dominen el análisis.
\end{enumerate}

\subsection{Planet frente a Sentinel-2}
Cada comparación utiliza 8,104 pares válidos de la misma parcela y fecha.
\begin{enumerate}[resume]
  \item \textbf{NDVI.} La correlación es 0.974, pero Planet es en promedio 0.083
  unidades mayor y el MAE es 0.083. Los puntos siguen una relación fuerte, aunque
  desplazada respecto a la igualdad; se requiere calibración antes de fusionar.

  \item \textbf{EVI.} La correlación es 0.958, la diferencia media es $-0.0013$
  y el MAE 0.0357. Es el índice con menor sesgo promedio, pero la dispersión
  impide afirmar igualdad observación por observación.

  \item \textbf{LAI.} La correlación baja a 0.910 y Planet excede a Sentinel-2
  por 1.062 unidades en promedio, con MAE de 1.062. El desplazamiento es
  sistemático y mucho mayor que en NDVI o EVI; ambos LAI no deben concatenarse
  directamente ni interpretarse en una escala común sin validación externa.
\end{enumerate}

\subsection{Saltos de NDVI Sentinel-2}
\begin{enumerate}[resume]
  \item El percentil 99 del cambio absoluto es 0.304 y produce 447 casos
  extremos en 169 parcelas. El salto máximo es 0.705. La nubosidad mediana entre
  estos casos es 0, pero el promedio es 25.0\% y el cuartil superior 62.1\%:
  algunas caídas fuertes sí coinciden con nubosidad alta, aunque también hay
  aumentos grandes con nubosidad cero. Por ejemplo, AGC\_150 sube 0.705 con
  nubosidad 0, mientras AGC\_154 cae 0.695 con 43.4\% de nubosidad. Por tanto,
  los saltos pueden combinar dinámica agrícola real, contaminación atmosférica
  y separación temporal. Deben revisarse individualmente junto con fecha previa,
  días transcurridos y calidad de ambas capturas; no deben eliminarse de forma
  automática.
\end{enumerate}

\section{Conclusiones que debe producir la revisión}
La lectura conjunta debe responder: si la nubosidad es comparable entre sensores;
si su ausencia o concentración en extremos compromete análisis posteriores; qué
parcelas tienen cobertura Landsat atípica; qué variables muestran sesgo entre
Planet y Sentinel-2; y qué saltos de NDVI requieren una revisión por parcela,
fecha, nubosidad e intervalo desde la observación anterior.

\end{document}

